\section{Introduction}
\label{sec:intro}

Tool-using LLM agents now read files, browse the web, run shell commands, and
send email on a user's behalf. This autonomy is also the attack surface: an
agent can be coaxed --- by a malicious user, by a poisoned web page, or by an
instruction smuggled into the output of a tool it called --- into actions that
leak credentials, destroy data, or escalate privilege. Unlike a single harmful
generation, an agent failure is a \emph{trajectory}: a sequence of
thought / action / observation steps in which an early benign-looking step
(read a config file) sets up a later harmful one (POST it to an external host).
The dangerous moment is a specific tool call, and once that call executes, the
damage is done.

\paragraph{The gap: nobody guards the mid-trajectory moment.}
Most deployed safeguards act at the wrong time or the wrong granularity.
\emph{Output filters} such as Llama Guard \citep{llamaguard} and WildGuard
\citep{wildguard} classify a model's final text; they are not built to gate the
next tool call mid-run. \emph{Post-hoc transcript judges} such as R-Judge
\citep{rjudge} score a completed trajectory --- useful for offline auditing, but
by construction they cannot prevent the action that already ran.
\emph{LLM-as-a-judge} per action would close the timing gap, but a frontier-scale
judge invoked before every tool call is too slow and too costly to sit in the
hot path of a live agent. What is missing is a guard that is (i) \emph{online}
--- it runs before each action and can veto it; (ii) \emph{cheap} --- small
enough to gate every step; and (iii) \emph{trajectory-aware} --- it scores the
proposed action in the context of everything that came before.

\paragraph{Contribution.}
We present \textbf{TraceGuard}, a sub-2B online trajectory monitor that fills
this gap. Our contributions are:
\begin{enumerate}[leftmargin=1.4em,itemsep=0.15em]
  \item \textbf{An online per-step verdict head.} A full fine-tune of
        Qwen3.5-2B renders the trajectory-so-far plus a proposed action and
        reads $P(\textsc{STOP})$ from the next-token logits over
        $\{\textsc{OK},\textsc{WARN},\textsc{STOP}\}$. The monitor's
        \texttt{check()} call returns a \textsc{STOP} verdict (plus a risk type
        and a natural-language reason) \emph{before} the action is dispatched,
        so the executor can halt it (Section~\ref{sec:method}).
  \item \textbf{Dense per-step verdict supervision.} Rather than a single
        end-of-trace label, we supervise a verdict slot after \emph{every}
        decision step --- \textsc{OK} on safe steps, \textsc{STOP} on the
        anomaly step --- which teaches the model ``do \emph{not} stop here'' as
        well as ``stop here'' and matches the deployment regime where the
        monitor decides at every step.
  \item \textbf{A framework-agnostic canonical schema.} We fold heterogeneous
        agent event formats into one canonical \texttt{Trajectory}
        (\texttt{\{steps: [\{role, thought, action, observation\}]\}}) with an
        ATBench-derived label taxonomy \citep{atbench}, so one model serves
        multiple frameworks (Claude Code, OpenAI-API ReAct, LangGraph)
        (Section~\ref{sec:data}).
  \item \textbf{A deployment-centric operating point.} We report
        \textbf{Prevention Rate at FSR$\le$1\%} (PR@FSR), which answers the
        question an operator actually asks: ``If I tolerate at most one false
        stop per hundred benign runs, what fraction of attacks do I prevent in
        time?'' (Section~\ref{sec:experiments}).
\end{enumerate}

\paragraph{Positioning: complement, not replacement.}
A well-aligned backend already refuses most agent attacks: in our red-team
probing, a production-aligned model resisted $0/65$ common attack prompts
across indirect injection, exfiltration, SSRF, and shell injection
(Section~\ref{sec:experiments}, F1). TraceGuard's value is for the substantial
fraction of production agents that run on \emph{weakly-aligned} backends, where
the same attacks succeed at high rates. We therefore frame agent attack success
as a two-dimensional phenomenon --- attacker goal $\times$ victim alignment ---
and position TraceGuard as a runtime layer that complements alignment for the
weakly-aligned regime rather than substituting for it.

On a preliminary 269-trace evaluation set, TraceGuard prevents
$\approx 89\%$ of attacks in time while false-stopping $\approx 2.2\%$ of
benign traces overall, and reaches PR@FSR=1\% $\approx 0.89$ at AUROC
$\approx 0.997$ --- far above a keyword baseline (PR@FSR=1\% $\approx 0.12$)
and an action-perplexity baseline ($\approx 0.08$). We are explicit about
limits: the evaluation set is small, results are from a single training run,
the residual false-stop rate is concentrated almost entirely in one benign
source, and latency / cross-backend transfer numbers are not yet measured
(Section~\ref{sec:limitations}).
